\chapter{Networking and HTTP}\label{ch:networking}

\index{networking}
\index{HTTP}
\index{TCP}
\index{TLS}

Most modern programs need to talk to the network. Whether you are fetching
data from a REST API, scraping a web page, or building a microservice,
networking is the backbone. Lattice provides networking at two levels: a
high-level HTTP client that handles the protocol details for you, and a
low-level TCP/TLS layer that gives you direct control over sockets. On top of
these primitives, the \texttt{http\_server} library offers an Express-style
framework for building web servers entirely in Lattice.

We will start at the top---making HTTP requests---and then work our way down
to raw sockets, TLS, and server construction.

% ---------------------------------------------------------------------------
\section{The http Module}\label{sec:http-module}
% ---------------------------------------------------------------------------

\index{http\_get}
\index{http\_post}
\index{http\_request}

The HTTP functions are available both as top-level built-ins
(\lstinline|http_get|, \lstinline|http_post|, \lstinline|http_request|) and
as methods on the \lstinline|http| module object. They handle URL parsing,
connection management, TLS negotiation for HTTPS, request formatting,
chunked transfer decoding, and response parsing---all behind a single
function call.

\subsection{Simple GET Requests}\label{sec:http-get}

\begin{lstlisting}[language=Lattice, caption={Making an HTTP GET request}]
let response = http_get("https://httpbin.org/get")

print(response["status"])   // 200
print(response["body"])     // JSON response body

// Response headers are a Map with lowercase keys
let headers = response["headers"]
print(headers["content-type"])  // "application/json"
\end{lstlisting}

\lstinline|http_get| takes a URL string and returns a \lstinline|Map| with
three keys:

\begin{center}
\begin{tabular}{llp{7.5cm}}
\toprule
\textbf{Key} & \textbf{Type} & \textbf{Description} \\
\midrule
\lstinline|"status"| & \lstinline|Int| & HTTP status code (200, 404, 500, etc.) \\
\lstinline|"body"| & \lstinline|String| & Response body as a string \\
\lstinline|"headers"| & \lstinline|Map| & Response headers (keys are lowercase) \\
\bottomrule
\end{tabular}
\end{center}

The URL must start with \lstinline|http://| or \lstinline|https://|. If the
scheme is \lstinline|https|, the HTTP layer automatically performs a TLS
handshake with certificate verification before sending the request---you do
not need to manage TLS yourself.

\begin{lstlisting}[language=Lattice, caption={Fetching and parsing JSON from an API}]
let response = http_get("https://api.example.com/users")

if response["status"] == 200 {
    let users = json_parse(response["body"])
    for user in users {
        print(user["name"] + " (" + user["email"] + ")")
    }
} else {
    print("Request failed with status " + to_string(response["status"]))
}
\end{lstlisting}

\begin{notebox}{HTTP and the Filesystem Chapter}
Notice how naturally HTTP combines with the data format functions from
\cref{ch:data-formats}. The response body is always a \lstinline|String|;
use \lstinline|json_parse|, \lstinline|yaml_parse|, or the other parsers to
turn it into structured data.
\end{notebox}

\subsection{POST Requests}\label{sec:http-post}

\begin{lstlisting}[language=Lattice, caption={Sending a POST request with a JSON body}]
let payload = Map::new()
payload["username"] = "alice"
payload["email"] = "alice@example.com"

let options = Map::new()
options["body"] = json_stringify(payload)
options["headers"] = Map::new()
options["headers"]["Content-Type"] = "application/json"

let response = http_post("https://api.example.com/users", options)
print(response["status"])  // 201
print(response["body"])
\end{lstlisting}

\lstinline|http_post| takes a URL and an optional options Map. The options
Map can include:

\begin{itemize}
    \item \lstinline|"body"| --- the request body as a String
    \item \lstinline|"headers"| --- a Map of additional request headers
    \item \lstinline|"timeout"| --- request timeout in milliseconds (default:
          30\,000)
\end{itemize}

\subsection{The General-Purpose http\_request}\label{sec:http-request}

For methods beyond GET and POST---PUT, PATCH, DELETE, HEAD---use the
general-purpose \lstinline|http_request|:

\begin{lstlisting}[language=Lattice, caption={Using http\_request for other methods}]
// PUT request to update a resource
let update_body = json_stringify({"name": "Alice Updated"})
let options = Map::new()
options["body"] = update_body
options["headers"] = Map::new()
options["headers"]["Content-Type"] = "application/json"

let response = http_request("PUT", "https://api.example.com/users/1", options)
print(response["status"])  // 200

// DELETE request
let del_response = http_request("DELETE", "https://api.example.com/users/1", {})
print(del_response["status"])  // 204
\end{lstlisting}

\lstinline|http_request| takes the HTTP method as a string, the URL, and
an optional options Map with the same keys as \lstinline|http_post|.

\begin{tipbox}{Timeouts}
All HTTP functions use a default timeout of 30 seconds. For slow APIs or
large downloads, pass a \lstinline|"timeout"| key in the options Map:
\begin{lstlisting}[language=Lattice]
let options = Map::new()
options["timeout"] = 60000  // 60 seconds
let response = http_get("https://slow-api.example.com/data")
\end{lstlisting}
A timeout of 0 means ``use the default'' (30 seconds), not ``no timeout.''
\end{tipbox}

\subsection{Error Handling in HTTP}\label{sec:http-errors}

HTTP functions raise a runtime error when something goes wrong at the
\emph{network} level: DNS resolution failure, connection refused, TLS
certificate validation failure, or timeout. HTTP-level errors (4xx, 5xx
status codes) are \emph{not} runtime errors---they are indicated by the
\lstinline|"status"| field in the response. This means you should always
check the status code:

\begin{lstlisting}[language=Lattice, caption={Robust HTTP error handling}]
fn fetch_user(user_id: Int) -> Map {
    let url = "https://api.example.com/users/" + to_string(user_id)

    let response = try {
        http_get(url)
    } catch error {
        print("Network error: " + to_string(error))
        return nil
    }

    if response["status"] != 200 {
        print("API error: " + to_string(response["status"]))
        return nil
    }

    return json_parse(response["body"])
}
\end{lstlisting}

\begin{warnbox}{HTTPS Certificate Verification}
Lattice's TLS layer verifies server certificates against the system's
trusted certificate store. If you connect to a server with an expired,
self-signed, or otherwise invalid certificate, the request will fail with a
TLS error. This is the correct, safe default---do not work around it in
production.
\end{warnbox}

% ---------------------------------------------------------------------------
\section{The net Module: TCP Networking}\label{sec:net-module}
% ---------------------------------------------------------------------------

\index{tcp\_listen}
\index{tcp\_accept}
\index{tcp\_connect}
\index{tcp\_read}
\index{tcp\_write}
\index{tcp\_close}
\index{tcp\_peer\_addr}
\index{tcp\_set\_timeout}

Beneath the HTTP layer, Lattice exposes raw TCP socket operations. These
functions work with integer \emph{socket descriptors}---opaque handles that
identify a connection. The net layer tracks which descriptors are valid
sockets, so passing a random integer will produce a clear error rather than
undefined behavior.

\subsection{Connecting to a Server}\label{sec:tcp-connect}

\begin{lstlisting}[language=Lattice, caption={Connecting to a TCP server}]
let fd = tcp_connect("example.com", 80)

// Send an HTTP request manually
let request = "GET / HTTP/1.1\r\nHost: example.com\r\nConnection: close\r\n\r\n"
tcp_write(fd, request)

// Read the response
let response = tcp_read(fd)
print(response)

tcp_close(fd)
\end{lstlisting}

\lstinline|tcp_connect| performs DNS resolution and establishes a TCP
connection. It returns an integer socket descriptor. \lstinline|tcp_write|
sends data over the connection, handling partial writes internally.
\lstinline|tcp_read| reads up to 8\,KB of available data and returns it as a
string. An empty string indicates end-of-file (the remote side closed the
connection).

\begin{lstlisting}[language=Lattice, caption={Reading all data from a TCP connection}]
fn read_all(fd: Int) -> String {
    flux result = ""
    loop {
        let chunk = tcp_read(fd)
        if chunk.len() == 0 {
            break  // EOF
        }
        result = result + chunk
    }
    return result
}
\end{lstlisting}

If you need to read an exact number of bytes, use \lstinline|tcp_read_bytes|:

\begin{lstlisting}[language=Lattice, caption={Reading an exact number of bytes}]
let fd = tcp_connect("example.com", 80)
tcp_write(fd, "GET / HTTP/1.1\r\nHost: example.com\r\nConnection: close\r\n\r\n")

// Read exactly 1024 bytes (or less if the connection closes)
let data = tcp_read_bytes(fd, 1024)
print("Received " + to_string(data.len()) + " bytes")

tcp_close(fd)
\end{lstlisting}

\subsection{Building a TCP Server}\label{sec:tcp-server}

To accept incoming connections, use \lstinline|tcp_listen| and
\lstinline|tcp_accept|:

\begin{lstlisting}[language=Lattice, caption={A minimal TCP echo server}]
let server_fd = tcp_listen("0.0.0.0", 9000)
print("Echo server listening on port 9000")

loop {
    let client_fd = tcp_accept(server_fd)
    let peer = tcp_peer_addr(client_fd)
    print("Connection from " + peer)

    // Read what the client sends
    let data = tcp_read(client_fd)

    // Echo it back
    tcp_write(client_fd, data)

    tcp_close(client_fd)
    print("Closed connection from " + peer)
}
\end{lstlisting}

\lstinline|tcp_listen| creates a server socket bound to the given address and
port, with \texttt{SO\_REUSEADDR} set so you can restart the server without
waiting for the socket to time out. \lstinline|tcp_accept| blocks until a
client connects, then returns a new socket descriptor for that connection.
\lstinline|tcp_peer_addr| returns the remote address as a
\lstinline|"ip:port"| string.

\begin{notebox}{Socket Timeouts}
Use \lstinline|tcp_set_timeout(fd, seconds)| to set read and write timeouts
on a socket. This is important for servers---without a timeout, a
misbehaving client that never sends data will block your server forever:
\begin{lstlisting}[language=Lattice]
let client_fd = tcp_accept(server_fd)
tcp_set_timeout(client_fd, 10)  // 10-second timeout
\end{lstlisting}
\end{notebox}

\subsection{TCP Function Reference}\label{sec:tcp-reference}

\begin{center}
\begin{tabular}{lp{8cm}}
\toprule
\textbf{Function} & \textbf{Description} \\
\midrule
\lstinline|tcp_listen(host, port)| & Bind and listen; returns server socket descriptor \\
\lstinline|tcp_accept(server_fd)| & Accept a connection; returns client socket descriptor \\
\lstinline|tcp_connect(host, port)| & Connect to a remote server; returns socket descriptor \\
\lstinline|tcp_read(fd)| & Read up to 8\,KB; returns String (empty on EOF) \\
\lstinline|tcp_read_bytes(fd, n)| & Read exactly \texttt{n} bytes; returns String \\
\lstinline|tcp_write(fd, data)| & Write all data; handles partial writes \\
\lstinline|tcp_close(fd)| & Close socket and release tracking \\
\lstinline|tcp_peer_addr(fd)| & Get remote \texttt{"ip:port"} as String \\
\lstinline|tcp_set_timeout(fd, secs)| & Set read/write timeout in seconds \\
\bottomrule
\end{tabular}
\end{center}

% ---------------------------------------------------------------------------
\section{TLS: Secure Connections}\label{sec:tls}
% ---------------------------------------------------------------------------

\index{tls\_connect}
\index{tls\_read}
\index{tls\_write}
\index{tls\_close}
\index{tls\_available}

For encrypted communication, Lattice provides TLS functions that mirror the
TCP API but add encryption and certificate verification. The TLS layer is
built on OpenSSL (or Schannel on Windows) and supports TLS 1.2 and above.

\subsection{Making a TLS Connection}

\begin{lstlisting}[language=Lattice, caption={Connecting with TLS}]
// Check if TLS support is available in this build
if !tls_available() {
    print("TLS not available -- build with OpenSSL support")
    return
}

let fd = tls_connect("example.com", 443)

let request = "GET / HTTP/1.1\r\nHost: example.com\r\nConnection: close\r\n\r\n"
tls_write(fd, request)

flux response = ""
loop {
    let chunk = tls_read(fd)
    if chunk.len() == 0 { break }
    response = response + chunk
}

print(response)
tls_close(fd)
\end{lstlisting}

\lstinline|tls_connect| performs a TCP connection, then negotiates a TLS
handshake on top of it. It sets the Server Name Indication (SNI) extension
and enables hostname verification, so the connection is secure by default.
The returned descriptor is the same integer type as TCP descriptors, but you
\emph{must} use the \lstinline|tls_*| functions to read and write---the data
must pass through the encryption layer.

\begin{warnbox}{Don't Mix TCP and TLS Functions}
A descriptor returned by \lstinline|tls_connect| must be used with
\lstinline|tls_read|, \lstinline|tls_write|, and \lstinline|tls_close|. If
you call \lstinline|tcp_read| on a TLS descriptor, you will read encrypted
gibberish. If you call \lstinline|tcp_close| instead of
\lstinline|tls_close|, the TLS session will not be properly shut down.
\end{warnbox}

\subsection{TLS Function Reference}\label{sec:tls-reference}

\begin{center}
\begin{tabular}{lp{8cm}}
\toprule
\textbf{Function} & \textbf{Description} \\
\midrule
\lstinline|tls_connect(host, port)| & TCP connect + TLS handshake + cert verification \\
\lstinline|tls_read(fd)| & Read up to 8\,KB of decrypted data \\
\lstinline|tls_read_bytes(fd, n)| & Read exactly \texttt{n} decrypted bytes \\
\lstinline|tls_write(fd, data)| & Encrypt and write all data \\
\lstinline|tls_close(fd)| & TLS shutdown + close socket \\
\lstinline|tls_available()| & Returns \lstinline|true| if TLS support was compiled in \\
\bottomrule
\end{tabular}
\end{center}

\begin{defnbox}{When to Use TLS Directly}
In practice, you rarely need the TLS functions directly. \lstinline|http_get|
and friends handle TLS automatically for \lstinline|https://| URLs. The TLS
layer is useful when you need to communicate over encrypted sockets with
non-HTTP protocols: database connections, message queues, custom binary
protocols, or connecting to services like Redis over TLS.
\end{defnbox}

% ---------------------------------------------------------------------------
\section{Building a Simple HTTP Server}\label{sec:http-server}
% ---------------------------------------------------------------------------

\index{http\_server}
\index{web server}

Now that we understand both the HTTP client and the TCP layer, let's build
something on the server side. Lattice ships with two approaches: you can
build a server from scratch using raw TCP (instructive), or use the
\texttt{http\_server} library for a more productive, framework-style
experience.

\subsection{A Server from Scratch}\label{sec:server-from-scratch}

The minimal approach constructs HTTP responses as raw strings:

\begin{lstlisting}[language=Lattice, caption={A minimal HTTP server from raw TCP}]
fn build_response(status: Int, status_text: String,
                  content_type: String, body: String) -> String {
    let header = "HTTP/1.1 " + to_string(status) + " " + status_text + "\r\n"
        + "Content-Type: " + content_type + "\r\n"
        + "Content-Length: " + to_string(body.len()) + "\r\n"
        + "Connection: close\r\n"
        + "\r\n"
    return header + body
}

let server = tcp_listen("0.0.0.0", 8080)
print("Listening on http://localhost:8080")

loop {
    let conn = tcp_accept(server)
    let raw = tcp_read(conn)

    if raw.len() > 0 {
        // Parse the request line: "GET /path HTTP/1.1"
        let lines = raw.split("\r\n")
        let parts = lines[0].split(" ")
        let method = parts[0]
        let path = if parts.len() > 1 { parts[1] } else { "/" }

        print(method + " " + path)

        let response = if path == "/" {
            build_response(200, "OK", "text/html",
                "<h1>Hello from Lattice!</h1>")
        } else if path == "/json" {
            build_response(200, "OK", "application/json",
                json_stringify({"message": "Hello!"}))
        } else {
            build_response(404, "Not Found", "text/plain",
                "404 Not Found: " + path)
        }

        tcp_write(conn, response)
    }
    tcp_close(conn)
}
\end{lstlisting}

This works, but it has obvious limitations: no routing, no middleware, no
cookie handling, and no request body parsing. The \texttt{http\_server}
library handles all of that.

\subsection{The http\_server Library}\label{sec:http-server-lib}

\index{http\_server library}

Import the library and create an application:

\begin{lstlisting}[language=Lattice, caption={Creating an HTTP server with the http\_server library}]
import "lib/http_server" as http

let app = http.new()
\end{lstlisting}

\subsubsection{Registering Routes}

Routes are registered by HTTP method:

\begin{lstlisting}[language=Lattice, caption={Registering route handlers}]
app.get("/", |req, res| {
    return http.response(200, "Hello, World!")
})

app.get("/about", |req, res| {
    return http.html(200, "<h1>About</h1><p>Built with Lattice.</p>")
})

app.post("/api/items", |req, res| {
    let body = req.get("body")
    print("Received: " + body)
    return http.json(201, {"status": "created"})
})

app.delete("/api/items", |req, res| {
    return http.json(200, {"status": "deleted"})
})
\end{lstlisting}

Route handlers receive a request Map and a response context Map. They must
return a response Map, which you build using the helper functions:

\begin{center}
\begin{tabular}{lp{8cm}}
\toprule
\textbf{Helper} & \textbf{Description} \\
\midrule
\lstinline|http.response(status, body)| & Plain text response \\
\lstinline|http.json(status, data)| & JSON response (auto-stringifies the data) \\
\lstinline|http.html(status, body)| & HTML response \\
\lstinline|http.redirect(url)| & 302 redirect (optional status parameter) \\
\lstinline|http.stream(status, chunks)| & Chunked transfer encoding response \\
\bottomrule
\end{tabular}
\end{center}

\subsubsection{The Request Map}

The request Map passed to handlers contains:

\begin{center}
\begin{tabular}{llp{7cm}}
\toprule
\textbf{Key} & \textbf{Type} & \textbf{Description} \\
\midrule
\lstinline|"method"| & \lstinline|String| & HTTP method (``GET'', ``POST'', etc.) \\
\lstinline|"path"| & \lstinline|String| & Request path without query string \\
\lstinline|"query"| & \lstinline|Map| & Parsed query parameters \\
\lstinline|"headers"| & \lstinline|Map| & Request headers \\
\lstinline|"body"| & \lstinline|String| & Request body \\
\lstinline|"cookies"| & \lstinline|Map| & Parsed cookies \\
\lstinline|"peer"| & \lstinline|String| & Remote address (\texttt{"ip:port"}) \\
\bottomrule
\end{tabular}
\end{center}

\begin{lstlisting}[language=Lattice, caption={Working with query parameters and cookies}]
app.get("/search", |req, res| {
    let query = req.get("query")
    let term = if query.has("q") { query.get("q") } else { "" }
    return http.json(200, {"results": [], "query": term})
})

app.get("/profile", |req, res| {
    let cookies = req.get("cookies")
    let session = if cookies.has("session") {
        cookies.get("session")
    } else {
        "anonymous"
    }
    let response = http.json(200, {"user": session})
    return http.set_cookie(response, "visited", "true", http.cookie_opts())
})
\end{lstlisting}

\subsubsection{Middleware}

Middleware functions run before route handlers. Each middleware receives the
request, response context, and a \lstinline|next| function to call the rest
of the pipeline:

\begin{lstlisting}[language=Lattice, caption={Adding middleware to a server}]
import "lib/http_server" as http

let app = http.new()

// Built-in logging middleware
app.use(http.logger())

// Built-in CORS middleware
app.use(http.cors())

// Built-in security headers
app.use(http.security_headers())

// Built-in JSON body parser
app.use(http.body_parser())

// Custom middleware: authentication
app.use(|req, res, next| {
    let headers = req.get("headers")
    if headers.has("Authorization") {
        req.set("authenticated", true)
    } else {
        req.set("authenticated", false)
    }
    return next(req, res)
})

app.get("/admin", |req, res| {
    if !req.get("authenticated") {
        return http.json(401, {"error": "Unauthorized"})
    }
    return http.json(200, {"message": "Welcome, admin"})
})

app.listen(8080)
\end{lstlisting}

The middleware pipeline executes in registration order. The \lstinline|logger|
middleware prints timing information for each request. The \lstinline|cors|
middleware handles preflight OPTIONS requests and sets the appropriate
\texttt{Access-Control-*} headers. The \lstinline|security_headers| middleware
adds \texttt{X-Content-Type-Options}, \texttt{X-Frame-Options}, and similar
defensive headers.

\begin{tipbox}{Custom CORS Configuration}
The \lstinline|cors()| middleware accepts an options Map:
\begin{lstlisting}[language=Lattice]
let cors_config = http.cors_opts()
cors_config["origin"] = "https://mysite.com"
cors_config["credentials"] = true
app.use(http.cors(cors_config))
\end{lstlisting}
\end{tipbox}

\subsubsection{Starting the Server}

\begin{lstlisting}[language=Lattice, caption={Starting the server}]
app.listen(8080)
// "Lattice HTTP server listening on http://localhost:8080"
\end{lstlisting}

The \lstinline|listen| call blocks forever, accepting connections in a loop.
Each connection is handled synchronously. For concurrent handling, you could
combine this with structured concurrency (\cref{ch:structured-concurrency}),
spawning a task for each accepted connection.

\subsection{A Complete REST API}\label{sec:rest-api}

Let's put it all together with a complete in-memory REST API:

\begin{lstlisting}[language=Lattice, caption={A complete REST API server}]
import "lib/http_server" as http

let app = http.new()
app.use(http.logger())
app.use(http.cors())
app.use(http.body_parser())

// In-memory data store
flux items = []
flux next_id = 1

// List all items
app.get("/api/items", |req, res| {
    return http.json(200, items)
})

// Create an item
app.post("/api/items", |req, res| {
    let body = req.get("body")
    let data = json_parse(body)
    data["id"] = next_id
    next_id += 1
    items.push(data)
    return http.json(201, data)
})

// Health check
app.get("/health", |req, res| {
    return http.json(200, {"status": "ok", "items": items.len()})
})

print("Starting API server...")
app.listen(3000)
\end{lstlisting}

% ---------------------------------------------------------------------------
\section{URL Encoding and Decoding}\label{sec:url-encoding}
% ---------------------------------------------------------------------------

\index{url\_encode}
\index{url\_decode}

URLs can only contain a limited set of ASCII characters. When you need to
include special characters---spaces, ampersands, Unicode text---you must
percent-encode them. Lattice provides \lstinline|url_encode| and
\lstinline|url_decode| for this purpose.

\begin{lstlisting}[language=Lattice, caption={URL encoding and decoding}]
let query = "name=Alice Bob&city=New York"
let encoded = url_encode(query)
print(encoded)  // "name%3DAlice%20Bob%26city%3DNew%20York"

let decoded = url_decode(encoded)
print(decoded)  // "name=Alice Bob&city=New York"
\end{lstlisting}

\lstinline|url_encode| percent-encodes every character that is not an
unreserved character (A--Z, a--z, 0--9, hyphen, underscore, period, tilde).
\lstinline|url_decode| reverses the process, converting \texttt{\%XX}
sequences back to the original bytes and converting \texttt{+} to spaces.

\begin{lstlisting}[language=Lattice, caption={Building a URL with query parameters}]
fn build_url(base: String, params: Map) -> String {
    let keys = params.keys()
    if keys.len() == 0 {
        return base
    }

    flux parts = []
    for key in keys {
        let encoded_key = url_encode(key)
        let encoded_val = url_encode(params[key])
        parts.push(encoded_key + "=" + encoded_val)
    }

    return base + "?" + parts.join("&")
}

let params = Map::new()
params["q"] = "lattice programming"
params["page"] = "1"
params["lang"] = "en"

let url = build_url("https://search.example.com/api", params)
print(url)
// https://search.example.com/api?q=lattice%20programming&page=1&lang=en
\end{lstlisting}

\begin{notebox}{Module vs. Top-Level Access}
\lstinline|url_encode| and \lstinline|url_decode| are available both as
top-level built-in functions and as methods on the \lstinline|crypto| module
(since they are conceptually about encoding/decoding). The top-level versions
are convenient for scripts; the module versions are better for explicit
imports in larger programs.
\end{notebox}

% ---------------------------------------------------------------------------
\section*{Exercises}\label{sec:networking-exercises}
% ---------------------------------------------------------------------------

\begin{enumerate}
    \item \textbf{HTTP Health Checker.} Write a program that takes a list of
    URLs from a configuration file, makes an HTTP GET request to each one,
    and reports whether the service is up (2xx status), degraded (5xx), or
    unreachable (network error). Use \lstinline|try|/\lstinline|catch| to
    handle connection failures gracefully.

    \item \textbf{Web Scraper.} Write a function that fetches a web page
    with \lstinline|http_get|, extracts all URLs from \texttt{href="..."}
    attributes using string operations (or the regex module from
    \cref{ch:time-crypto-os}), and returns them as an array. Bonus: follow
    the links one level deep and report the status of each.

    \item \textbf{TCP Chat Server.} Extend the echo server from
    \cref{sec:tcp-server} into a simple chat server. When a client connects,
    ask for their name, then broadcast each message to all connected clients.
    You will need to use structured concurrency (\cref{ch:structured-concurrency})
    to handle multiple clients simultaneously.

    \item \textbf{REST API with Persistence.} Extend the REST API from
    \cref{sec:rest-api} to persist its data to a JSON file. Load the file
    on startup and save after each modification. Use the safe-write pattern
    from \cref{sec:safe-pipeline} in the filesystem chapter.

    \item \textbf{HTTP File Server.} Build a server that serves static
    files from a directory. Map the URL path to a file path, read the file,
    determine the content type from the extension, and send it back. Handle
    directory traversal attacks by checking that the resolved path stays
    within the root directory (use \lstinline|realpath|).
\end{enumerate}

% ---------------------------------------------------------------------------
\section*{What's Next}\label{sec:networking-next}
% ---------------------------------------------------------------------------

We have now covered the major I/O systems in Lattice: filesystems, data
formats, and networking. In the final chapter of this Part, we will survey the
remaining standard library modules---time and date handling, cryptographic
functions, UUIDs, operating system interaction, mathematics, and regular
expressions---giving you the complete toolkit for building real-world
applications.
